\documentclass[11pt,a4paper]{article}
\usepackage[T1]{fontenc}
\usepackage[utf8]{inputenc}
\usepackage{amsmath,amssymb,xcolor}
\usepackage[margin=25mm]{geometry}
\usepackage[colorlinks=true,urlcolor=blue]{hyperref}
\definecolor{nexorablue}{HTML}{2459B5}
\setlength{\parindent}{0pt}
\setlength{\parskip}{7pt}
\setlength{\emergencystretch}{2em}
\renewcommand{\familydefault}{\sfdefault}
\begin{document}
\makeatletter
\def\@oddfoot{\footnotesize Download for free at \url{https://openstax.org/details/books/university-physics-volume-2}.\hfil\thepage}
\let\@evenfoot\@oddfoot
\makeatother
{\small\bfseries\color{nexorablue} NEXORA CPGE}\par
{\small MP / OST-SPE-13-01 / Corrigé}\par
{\LARGE\bfseries Bobine carrée dans un champ décroissant}\par
\vspace{0.25em}\hrule\vspace{0.6em}

\section*{Flux lié}
Pour la normale choisie, chaque spire reçoit $\Phi=B\ell^2$. La dérivée du flux lié est $N\ell^2\dot B$.

\section*{Force électromotrice et courant}
D’après Faraday, la force électromotrice algébrique est positive, car $\dot B<0$.
\[ |e|=200(0{,}25)^2(0{,}040)\,\mathrm V=0{,}50\,\mathrm V,\qquad |i|=\frac{0{,}50}{5{,}0}\,\mathrm A=0{,}10\,\mathrm A. \]

\section*{Contrôle du sens}
Le courant positif crée un champ dans le même sens que la normale. Il compense partiellement la diminution du flux excitateur.

\bigskip
{\footnotesize Adaptation française par Nexora CPGE. Source : Samuel J. Ling, Jeff Sanny et William Moebs (OpenStax, Rice University),
\href{https://d3bxy9euw4e147.cloudfront.net/oscms-prod/media/documents/UniversityPhysicsVolume2-OP.pdf\#page=595}{\textit{University Physics Volume 2, édition 2018}}, § 13.1, exemple 13.1, p. 585.
\href{https://creativecommons.org/licenses/by/4.0/}{CC BY 4.0}.}
\end{document}
